\documentclass[a4paper,12pt]{letter}

\usepackage{amsmath}
\usepackage{amssymb}
\usepackage{nopageno}
\usepackage{amsmath}
\usepackage{enumitem}
% \usepackage{mathtools}
\usepackage{eqnarray,amsmath}
\usepackage[a4paper, total={7.5in, 11.25in}]{geometry}
%\usepackage[margin=2.5cm]{geometry}
\usepackage[utf8]{inputenc}
\usepackage[T1]{fontenc}
\usepackage{lmodern}
%\usepackage{amsmath, amssymb}
\usepackage{setspace}
\setstretch{1.2}

\begin{document}

\begin{center}
 (KTXFI2EBNF) Physics II.  Exercises\\ \smallskip
November~3,~2025
 \end{center}

\begin{enumerate}

\item A light source has a power output of $75\,W$.

\begin{enumerate}
    \item[a)] What is the frequency of the light if the entire energy of the source is emitted as light of wavelength $ \lambda = 600\,nm $?
    \item[b)] How many photons are emitted by the light source per second?
\end{enumerate}

\item How many photons per second are emitted by a $ 7.5\,mW $ carbon dioxide laser if its wavelength is $ 10.6\,\mu m$?

\item A laser pulse of energy $ 20 \,J$ has a duration of $ 5 \cdot 10^{-7} \,$.  The light, of wavelength $ \lambda = 580\,nm $, strikes a metal surface.

\begin{enumerate}
    \item What is the power of the laser?
    \item How many photons strike the metal surface?
    \item What is the maximum velocity of the electrons ejected from the metal by the light, if the work function is $ 3 \cdot 10^{-19} \,J$?     (Electron mass: $ m_{e} = 9.11 \cdot 10^{-31} \,kg$)
\end{enumerate}

\item What is the energy of a photon of X-ray radiation with a wavelength of $ 10^{-10}\,m$, expressed in joules and in electronvolts?  What are the photon's momentum and mass?

\item A metal surface is illuminated with light of wavelength $ \lambda_{1} = 492\,nm $.  The maximum kinetic energy of the emitted electrons is $ W_{kin1} = 7.9 \cdot 10^{-19} \,J$.  When the wavelength of the light is changed to $ \lambda_{2} = 579\,nm $, the kinetic energy decreases to $ W_{kin2} =3.8 \cdot 10^{-20} \,J$.  Based on these data, determine the Planck constant and the work function characteristic of the metal!

\item When light of a certain frequency is incident on a metal surface, the measured stopping potential of the emitted electrons is $ U_{f1} = 3.19\,V$.  For light of half that frequency, the stopping potential is $ U_{f2} = 0.625 \,V $.   Determine the work function and the wavelength of the light!
\end{enumerate}

  
\bigskip

\rightline{Dr.~Gabor~FACSKO}
\rightline{\textit{facsko.gabor@uni-obuda.hu}}

\vfill

\end{document}
